\section{Proxy vs Exact: Three Scenarios}
  \label{sec:comparison}

  The comparison between $\hat{\delta}_{\mathrm{exact}}$ and the proxy value
  $\hat{\delta}_{\mathrm{cap}}\approx 3.4$ from O11~\cite{Beau2026a15} admits three
  qualitatively distinct outcomes, each with a different structural interpretation.
  We analyse them in turn, then identify the structural mechanism that would explain
  the discrepancy or agreement in each case.

  \subsection{Case A: $\hat{\delta}_{\mathrm{exact}}\approx\hat{\delta}_{\mathrm{cap}}$}

    If the exact and proxy exponents agree to within the uncertainty bands reported in
    Table~\ref{tab:sensitivity}, the central coordinate $\gamma$ is dynamically
    negligible at the asymptotic scale probed by these primes.
    The interpretation is that the capacity decay is governed by the translational and
    phase structure encoded in the $(a,b)$-coordinates alone, and the metaplectic
    twist $\psi_c(\gamma)$ averages out across BFS shells without affecting the
    leading-order rate.

    This outcome would retroactively validate O11: the proxy observable is a faithful
    instrument for $\delta$, and the result $\hat{\delta}_{\mathrm{cap}}\approx 3.4$
    already constitutes the first reliable measurement.
    It would also suggest that the asymptotic regime is already well-established at $q=307$
    and that the proxy-level underestimate identified in O10 was of algorithmic rather
    than structural origin.

  \subsection{Case B: $\hat{\delta}_{\mathrm{exact}}>\hat{\delta}_{\mathrm{cap}}$}

    If the exact exponent exceeds the proxy value, the central coordinate $\gamma$
    accelerates the capacity decay.
    The mechanism is the following.
    In the exact projection, the phase $\psi_c(\gamma+bx)$ introduces coherent
    interference between path contributions from different shells: paths that land
    on group elements with different $\gamma$-coordinates but the same $(a,b)$-projection
    are now distinguishable in $H_c$, contributing genuinely new directions to the span.
    This increases the incremental rank $\Delta r_n^{(c)}$ at early depths (when
    $\gamma$-diversity is high), which raises $\Sigma_n^{(c)}$ at small $n$, but
    the span saturates faster because the effective ambient dimension is $q$ rather than
    $q^2$ — leading to a steeper decay at intermediate depths.

    This outcome implies that O11 systematically underestimated $\delta$, and that the
    phenomenological window for $\beta^*$ must be recomputed using $\hat{\delta}_{\mathrm{exact}}$
    as input to the O3--O7 structural relation.

  \subsection{Case C: $\hat{\delta}_{\mathrm{exact}}<\hat{\delta}_{\mathrm{cap}}$}

    If the exact exponent falls below the proxy value, the central coordinate $\gamma$
    decelerates the capacity decay.
    The mechanism here is the opposite of Case B: the metaplectic phase $\psi_c(\gamma+bx)$
    introduces destructive interference between path contributions, reducing the
    effective incremental rank at each shell.
    This would mean that paths which the proxy counted as independent are in fact
    correlated in the exact representation via the central extension, so the span
    grows more slowly and $\delta$ is smaller.

    This outcome does not invalidate the programme but imposes a downward revision of
    $\hat{\delta}$ and a corresponding widening of the constraint on $\beta^*$.
    It would also raise a deeper question: whether the slower decay reflects a genuine
    property of the physical dynamics or an artefact of insufficient $q$ (i.e.\ the
    asymptotic regime not yet reached).
    The convergence of $\hat{\delta}_{\mathrm{exact}}$ with $q$ in Table~\ref{tab:delta}
    is the primary diagnostic for distinguishing these two interpretations.

  \subsection{Structural explanation via the metaplectic phase}
    \label{sec:structural-explanation}

    All three cases can be understood through a single structural parameter: the
    effective coherence length $\ell_\gamma$ of the central phase across a BFS shell.
    Formally, define
    \begin{equation}
      \ell_\gamma(n) = \left|\frac{1}{|S_n|}\sum_{g\in S_n}
                         e^{2\pi i c\,\gamma_g/q}\right|,
      \label{eq:coherence}
    \end{equation}
    where $\gamma_g$ is the central coordinate of $g\in S_n$.
    If $\ell_\gamma(n)\approx 1$ (high coherence), the central phase acts as a nearly
    constant multiplicative factor and the dynamics reduces to the proxy case (Case A).
    If $\ell_\gamma(n)\ll 1$ (low coherence, $\gamma_g$ uniformly distributed), the central
    phase randomises path contributions and modifies $\delta$ in either direction
    depending on whether the randomisation enriches or impoverishes the span.

    The regime of $\ell_\gamma$ as a function of $n$ can be extracted from the numerical
    data as a byproduct of the span-tracking algorithm, providing an independent diagnostic
    of the mechanism responsible for the observed discrepancy or agreement.

  \subsection{Dependence on $q$: proxy artefact or asymptotic convergence?}

    A fourth scenario is possible: $\hat{\delta}_{\mathrm{exact}}$ may depend on $q$ at
    the primes tested, with convergence occurring only at $q\gg 401$.
    If the exponents at $q\in\{211,307,401\}$ differ by more than $\sigma_{\mathrm{seed}}$
    without forming a clear limit, neither Case A, B, nor C applies cleanly.
    This scenario is identified in advance so that Table~\ref{tab:delta} can be read with
    it in mind: a monotone trend in $\hat{\delta}_{\mathrm{exact}}$ with $q$ that has not
    yet flattened would be evidence for an asymptotic artefact rather than a structural
    effect of $\gamma$.
